\section{Materials and methods}
\noindent\textit{Draft status. This section distinguishes documented procedures and counts from the protocol for remaining work. Corpus-wide full-text eligibility and a final included-study count have not been established.}

\subsection{Review design and analytical scope}
The review follows the information pathway from understanding existing buildings to assessing the environmental consequences of managing their materials. Five analytical stages organise this pathway: building observation and attribute inference (S1), material inventory (S2), stock changes and material flows over time (S3), recovery and subsequent-use pathways (S4), and environmental assessment and decision support (S5). These stages guide the comparison of integrated workflows reported in the literature. Their boundaries and connections are examined against the methods used in the studies.

Publications are screened individually. Within each publication, analysis focuses on how information passes from one model or stage to another. A paper may inform several stages but is counted only once in the publication total. Where several publications describe the same underlying study, their relationship must be checked before counting them as independent studies.

\subsection{Search families and information sources}
The searches cover Scopus, Web of Science Core Collection and IEEE Xplore. Full search expressions, search dates, field scopes and export logs are retained in the supplementary search records.

Base denotes searches on building material assessment. U denotes uncertainty, and DT denotes digital twin. Four search families are distinguished: Base; Base plus U; Base plus DT; and Base plus DT plus U. General DT and DT-plus-U searches additionally examine definitions and uncertainty mechanisms beyond the building-material domain. A publication may appear in more than one family. Some DT searches also use broader building and observation terms than the original Base searches. Their totals are therefore reported separately rather than interpreted as directly comparable proportions of a single common corpus. The supplementary protocol preserves the exact expressions needed to identify these differences. The additional integrated-framework search has no imposed year or publisher restriction; its selection log remains to be completed.

Scopus used title, abstract and keyword fields; Web of Science used Topic searches, including Keywords Plus; IEEE used the recorded search scope and native export format. These field scopes are not identical. The DT-stage protocol records no intentional year, language, document-type or open-access filter. Restrictions in earlier search versions must be taken from their own logs. A displayed result total is distinguished from a successfully exported record set. In particular, general DT annual facets do not constitute a complete record-level export.

\subsection{Search concepts and historical query differences}
The stage-specific concepts comprise spatial observation and building-attribute inference (S1); material stocks, inventories, intensities and archetypes (S2); lifetimes, demolition, stock turnover and temporal outflows (S3); recovery, reuse, recycling and circular routing (S4); and life-cycle assessment, embodied emissions, environmental impacts and decision support (S5). Uncertainty terms cover uncertainty, probabilistic and stochastic methods, Bayesian approaches, Monte Carlo simulation, sensitivity, variability, calibration and interval estimates. Exact term lists vary across archived versions and are preserved in Supplement S1.

Repository recovery established that S1--S3 Scopus expressions and split IEEE variants are available, together with executed WoS S1--S2 expressions. IEEE subquery results require consolidation before publication-level counting. Historical Scopus queries contain explicit NOT blocks; notably, the S3 Base and uncertainty versions differ in their exclusions. Thus, Base plus U cannot universally be interpreted as a strict subset of Base. Broad exclusion blocks can also suppress relevant multidisciplinary studies and are a limitation of retrieval.

The WoS S3 log and archived exports require reconciliation: the log labels S3-2 as not yet executed, although an S3-2 export exists. The executed query filenames referenced by that log were not found in the inspected repository tree. This discrepancy is retained in the supplement rather than resolved by inventing a historical expression.

\subsection{Eligibility by evidence stream}
Eligibility is assessed separately for the three synthesis streams. Stage--uncertainty evidence must identify a material-chain quantity and explain its uncertainty or a transformation affecting it. Stage--DT evidence must describe a representation, model or connection relevant to the building-material workflow. DT--uncertainty evidence may originate outside the building domain when the mechanism is sufficiently described to assess transferability; such evidence is not proof of implementation in the circular built environment.

An interface requires an identifiable input, transformation and downstream output. Mentioning two stages is insufficient. Unrelated publications are excluded, whereas useful indirect methods and secondary reviews are identified separately from direct domain evidence. Reviews do not count as independent replications of their cited primary studies. Insufficient information at title/abstract screening requires further assessment rather than an inferred full-text conclusion.

\subsection{Record management, screening and retrieval}
Native exports were preserved and bibliographic fields harmonised while retaining source and query membership. Exact normalised DOIs were used for conservative merging; native database identifiers preserved identity where DOIs were absent. Similar titles and conflicting DOIs were flagged for adjudication rather than automatically equated. Duplicate decisions and bibliographic corrections were retained in audit files. Deduplication status differs across families and is reported explicitly.

Title and abstract screening distinguished direct material-chain relevance, complementary methods, secondary evidence and unrelated records. Operational-only energy studies, generic project management and isolated material tests without the required assessment connection were distinguished from core material-flow evidence. An uncertainty or DT keyword alone did not establish eligibility. Incomplete information was retained as unresolved where a full text was necessary. Background and transferable-method records were retained separately rather than represented as core inclusions.

The selection process is organised using PRISMA 2020 reporting terminology (Page et al., 2021, DOI: 10.1136/bmj.n71). The flow diagrams distinguish the following steps:
\begin{enumerate}
\item Identification: records are collected from each database and additional search route, retaining their source identifiers.
\item Removal before screening: repeated records are consolidated and the reasons for any other removals are recorded.
\item Title and abstract screening: records are checked against the review scope; clearly unrelated records are separated and uncertain cases proceed to further assessment.
\item Full-text retrieval: reports selected for closer assessment are sought, with unavailable reports recorded separately.
\item Full-text eligibility: retrieved reports are assessed against the relevant evidence-stream criteria, with a specific reason recorded for each exclusion.
\item Inclusion and synthesis: eligible reports are assigned to evidence streams, retaining overlaps and distinguishing publication counts from independent-study counts.
\end{enumerate}
This sequence defines the reporting structure; its inclusion here does not imply completion of every step for every search family. The diagrams show verified counts and explicitly mark unreconciled stages. Downloaded files, full-text appraisals and final inclusions are distinct. Detailed screening provenance is retained in the audit materials and must be reconciled for final reporting.

\subsection{Digital twin definitions and data extraction}
Definitions are treated as evidence to be compared rather than resolved by the presence of the term digital twin. The extraction protocol separately records the authors' definition and cited definition source, the physical referent, digital representation, data connection, update mechanism and demonstrated analytical or decision function. This operational description draws on the characterisation and definition literature already cited in the introduction; it is not asserted to be a universally accepted threshold.

Manual versus automated exchange, one-way versus two-way exchange, periodic versus event-driven updates, and prediction versus feedback control are recorded separately. A BIM, GIS or three-dimensional model is not automatically coded as an operational twin. Conversely, absence of a reported feature is coded as not reported, not as proof that the feature is absent. Conceptual architectures, prototypes and operational deployments are distinguished. Domain studies that provide useful information mechanisms without satisfying an operational DT description remain identifiable as complementary evidence.

For uncertainty, the protocol records the affected quantity, source, representation, quantification method, dependence assumptions, transformation and downstream consequence. Spatial variability, epistemic uncertainty and alternative scenarios are not collapsed into one probabilistic category. Interface coding distinguishes documented transfer, documented representation change or loss, and insufficient reporting. Any claim of reduced uncertainty requires a stated comparison basis. Spatial extent and grain, temporal validity and horizon, and stakeholder information requirements are recorded as contextual conditions. Each substantive mechanism claim should be linked to a page, section, equation or figure; missing locators remain explicit. These fields extend existing appraisals and are not yet complete across all retrieved records.

\subsection{Full-text appraisal and evidence strength}
For each report, extraction records its version, source, eligibility decision and reason, together with page, section, equation, table or figure locators for substantive claims. Methods, results and limitations must be read together to distinguish a proposed feature from a demonstrated capability. Existing selected-paper labels and downloaded files do not substitute for this eligibility decision.

Appraisal addresses data provenance, model assumptions, dependency handling, validation and reproducibility. No universal numerical quality score is imposed across dissimilar study designs. Instead, these dimensions qualify the strength and scope of each synthesis claim. Missing reporting is distinguished from demonstrated absence, and uncertainty reduction requires an explicit comparison basis. The final extraction register must document these dimensions for every included report.

\subsection{Comparative analysis and evidence synthesis}
Existing relational analyses aggregate coded working appraisals into stage--uncertainty, stage--DT-function and DT--uncertainty-mechanism relationships. The denominator is stated for each analysis. Multiple memberships are preserved, and counts across stages or streams are not summed as independent publications. Frequencies indicate representation in the assessed set, not effectiveness or absence from the wider field. Retrieval-only records are excluded from full-text mechanism claims.

The intended final synthesis links an observed interface finding to an information requirement, a reported DT capability and its conditions of transfer. Directly demonstrated relationships are distinguished from cross-study interpretations and research proposals. A conceptual architecture is an organising result of the review, not an implemented or validated system. Literature silence is separated from a documented limitation. No meta-analysis or pooled effectiveness estimate is planned because the current evidence combines heterogeneous methods, contexts and study types.


\subsection{Analytical denominators and reproducibility}
Temporal trends based on summed database counts represent database occurrences, not unique publications. Unique-publication trends require record-level consolidation. Full-text content analyses use only reports with the required assessed evidence; metadata-only records do not support mechanism claims. Multiple stage memberships are retained, so sums across categories may exceed the number of publications.

Graph edges and matrix cells must be traceable to coded publication identifiers. Co-occurrence is not evidence of causation or effectiveness, and visual contrast alone does not establish statistical significance. Search expressions, source memberships, duplicate resolutions, eligibility reasons and extraction locators provide the audit trail from retrieval to synthesis. Historical noise-diagnostic scripts support topic diagnostics but do not establish completed full-text review.

\subsection{Reporting completeness and limitations}
The supplementary audit separates search occurrences, deduplicated publications, screening candidates, retrieved reports, full-text appraisals and final inclusions. The current diagrams are progress snapshots because full-text eligibility and some search-family joins remain incomplete. Counts are not pooled across overlapping families without a record-level reconciliation. Remaining work includes early-search provenance checks, unresolved duplicates, full-text decisions and supplemental-search accounting. The conceptual synthesis will be bounded to the evidence actually assessed.

